%% Lecture 22 -- Dipole Radiation Power and the Energy-Momentum Tensor
\section{Lecture 22 -- Radiated Power and the Energy-Momentum Tensor}\label{sec:lec22}

\begin{center}
\begin{tabular}{ll}
  \textbf{Date:}\quad 21/06/2026 & \\
\end{tabular}
\end{center}
\hrule\vspace{1.5em}

\subsection*{Overview}

This lecture closes the radiation calculation started in Lecture~21 and then uses it
as motivation for a more systematic discussion of energy and momentum in field theory.
The first half begins with the far-zone electric and magnetic fields of an accelerated
electric dipole. Because the radiation fields fall as $1/r$, their Poynting flux falls
as $1/r^2$, exactly compensating the area $4\pi r^2$ of a large sphere. This gives a
finite total power carried to infinity.

The second half asks where this energy accounting comes from in the action language.
Just as ordinary mechanics derives conserved energy and momentum from spacetime
translation symmetry, a classical field theory derives a conserved tensor
$T^\nu{}_\mu$ from the absence of explicit $x^\mu$ dependence in the Lagrangian
density. The lecture derives the Euler--Lagrange equations for a general field
$\psi(x)$ and then constructs the canonical energy-momentum tensor as the object whose
four divergence vanishes.

% ==============================================================================
\subsection{Radiation Fields of an Electric Dipole}
\label{subsec:lec22:radiation-fields}
% ==============================================================================

\subsubsection*{Physical Motivation}
Lecture~21 produced the far-zone fields of a localized, non-relativistic source whose
leading time-dependent multipole is the electric dipole moment $\bvec{p}(t)$. The
source size is $L$, the observation point is at radius $r\gg L$, and all source
quantities are evaluated at the retarded time
\begin{equation}\label{eq:lec22:retarded-time}
  t_r \equiv t-\frac{r}{c}.
\end{equation}
Only terms in the fields that decay as $1/r$ can transport a nonzero amount of energy
to infinity. Near-field terms such as the Coulomb field decay faster and therefore do
not contribute to the total radiated power at $r\to\infty$.

\dfn[def:lec22:radiation-zone]{Radiation Zone Fields}{
For electric dipole radiation,
\begin{align}
  \bB(t,\bvec{x})
  &= \frac{\ddot{\bvec{p}}(t_r)\times\uvec{r}}{c^2 r},
  \label{eq:lec22:B-rad}\\
  \bE(t,\bvec{x})
  &= \bB(t,\bvec{x})\times\uvec{r}
   = \frac{\qty[\ddot{\bvec{p}}(t_r)\times\uvec{r}]\times\uvec{r}}{c^2 r}.
  \label{eq:lec22:E-rad}
\end{align}
Here $\uvec{r}$ points from the source to the observer.
}

\textit{Significance:} Radiation is controlled by $\ddot{\bvec{p}}$, so a dipole
radiates only when the charges accelerate.

\subsubsection*{Geometry of the Fields}
Let $\theta$ be the angle between $\ddot{\bvec{p}}(t_r)$ and $\uvec{r}$. From
\eqref{eq:lec22:B-rad},
\begin{equation}\label{eq:lec22:B-magnitude}
  |\bB|=\frac{|\ddot{\bvec{p}}(t_r)|\sin\theta}{c^2 r}.
\end{equation}
The electric field is the transverse part of $-\ddot{\bvec{p}}/(c^2r)$:
\begin{align}
  \bE
  &= \frac{\qty[\ddot{\bvec{p}}\times\uvec{r}]\times\uvec{r}}{c^2 r} \notag\\
  &= \frac{\uvec{r}\qty(\uvec{r}\cdot\ddot{\bvec{p}})-\ddot{\bvec{p}}}{c^2 r}
   = -\frac{\ddot{\bvec{p}}_{\perp}}{c^2 r},
  \label{eq:lec22:E-transverse}
\end{align}
where $\ddot{\bvec{p}}_{\perp}$ is the component perpendicular to $\uvec{r}$. Therefore
\begin{equation}\label{eq:lec22:E-B-triad}
  \boxed{\bE\perp\bB,\qquad \bE\perp\uvec{r},\qquad \bB\perp\uvec{r},
  \qquad |\bE|=|\bB|.}
\end{equation}
\textit{Significance:} Far from the source, the spherical wave locally has the same
transverse structure as a plane electromagnetic wave.

% ==============================================================================
\subsection{Poynting Vector and Angular Power Distribution}
\label{subsec:lec22:poynting}
% ==============================================================================

\subsubsection*{Physical Motivation}
The radiation fields are not just oscillating fields; they carry energy away from the
source. In Gaussian units the energy flux density is the Poynting vector
\begin{equation}\label{eq:lec22:poynting-definition}
  \boxed{\bvec{S}=\frac{c}{4\pi}\,\bE\times\bB.}
\end{equation}
\textit{Significance:} $\bvec{S}\cdot\uvec{n}$ is the electromagnetic energy crossing a
unit area per unit time through a surface with normal $\uvec{n}$.

\subsubsection*{Derivation}
Using the radiation-zone triad \eqref{eq:lec22:E-B-triad}, the cross product points
radially outward:
\begin{equation}\label{eq:lec22:S-direction}
  \bE\times\bB = |\bE|\,|\bB|\,\uvec{r}=|\bB|^2\uvec{r}.
\end{equation}
Substituting \eqref{eq:lec22:B-magnitude},
\begin{align}
  \bvec{S}
  &= \frac{c}{4\pi}\,
     \frac{|\ddot{\bvec{p}}(t_r)|^2\sin^2\theta}{c^4 r^2}\,\uvec{r} \notag\\
  &= \boxed{
     \frac{|\ddot{\bvec{p}}(t_r)|^2\sin^2\theta}{4\pi c^3 r^2}\,\uvec{r}}.
  \label{eq:lec22:S-dipole}
\end{align}
\textit{Significance:} The $1/r^2$ falloff of the flux is exactly what is needed for a
finite power through a sphere of radius $r$.

The power per solid angle is obtained by multiplying by the area element
$r^2\dd{\Omega}$:
\begin{equation}\label{eq:lec22:angular-power}
  \boxed{\frac{dP}{d\Omega}
  = r^2\,\bvec{S}\cdot\uvec{r}
  = \frac{|\ddot{\bvec{p}}(t_r)|^2}{4\pi c^3}\sin^2\theta.}
\end{equation}
\textit{Significance:} Electric dipole radiation has a doughnut-shaped angular pattern:
zero on the dipole axis and maximal in the plane perpendicular to the dipole.

% ==============================================================================
\subsection{Total Radiated Power}
\label{subsec:lec22:total-power}
% ==============================================================================

\subsubsection*{Mathematical Setup}
The total power crossing a sphere of radius $r$ is
\begin{equation}\label{eq:lec22:power-surface}
  P = \oint_{S_r} \bvec{S}\cdot d\bvec{\sigma},
  \qquad d\bvec{\sigma}=r^2\sin\theta\dd{\theta}\dd{\phi}\,\uvec{r}.
\end{equation}
The limit $r\to\infty$ isolates energy that genuinely escapes the source instead of
merely sloshing in the near field.

\subsubsection*{Complete Angular Integral}
Insert \eqref{eq:lec22:S-dipole}:
\begin{align}
  P
  &= \int_0^{2\pi}\!\dd{\phi}\int_0^\pi\!\dd{\theta}\,
     \frac{|\ddot{\bvec{p}}(t_r)|^2\sin^2\theta}{4\pi c^3 r^2}
     r^2\sin\theta \notag\\
  &= \frac{|\ddot{\bvec{p}}(t_r)|^2}{4\pi c^3}
     \qty(\int_0^{2\pi}d\phi)
     \qty(\int_0^\pi \sin^3\theta\dd{\theta}).
  \label{eq:lec22:power-before-integral}
\end{align}
For the remaining integral, set $u=\cos\theta$, $du=-\sin\theta\dd{\theta}$:
\begin{align}
  \int_0^\pi \sin^3\theta\dd{\theta}
  &= \int_0^\pi (1-\cos^2\theta)\sin\theta\dd{\theta} \notag\\
  &= \int_{1}^{-1} (1-u^2)(-du)
   = \int_{-1}^{1}(1-u^2)\dd{u} \notag\\
  &= \qty[u-\frac{u^3}{3}]_{-1}^{1}
   = \qty(1-\frac13)-\qty(-1+\frac13)
   = \frac{4}{3}.
  \label{eq:lec22:sin3-integral}
\end{align}
Thus
\begin{align}
  P
  &= \frac{|\ddot{\bvec{p}}(t_r)|^2}{4\pi c^3}
     \qty(2\pi)\qty(\frac{4}{3}) \notag\\
  &= \boxed{\frac{2}{3c^3}\,|\ddot{\bvec{p}}(t_r)|^2.}
  \label{eq:lec22:dipole-power}
\end{align}
\textit{Significance:} This is the electric dipole radiation power in Gaussian units.
It is positive and finite at infinity, so the source must lose mechanical energy or be
driven by external work.

\nt{Only the $1/r$ pieces of $\bE$ and $\bB$ survive in the radiated power. If a field
piece decays as $1/r^2$, then its flux scales at most as $r^2/r^3$ or $r^2/r^4$ after
forming $\bE\times\bB$, and therefore vanishes as the sphere is sent to infinity.}

% ==============================================================================
\subsection{Why Radiation Does Not Violate Energy Conservation}
\label{subsec:lec22:energy-conservation}
% ==============================================================================

\subsubsection*{Physical Motivation}
A student asked why there is no contradiction between a positive outward flux at
infinity and conservation of energy. The answer is that radiation is not created from
nothing: accelerated charges require forces, and those forces do work. The work done
on the charges can be converted into electromagnetic energy that propagates away.

\begin{remark}
The classical atom is the standard warning example. A point electron in circular orbit
around a nucleus is constantly accelerating. Classical electrodynamics then predicts
dipole radiation, so the electron loses mechanical energy and spirals inward. The fact
that ordinary matter is stable is one of the historical signs that atomic physics
cannot be purely classical.
\end{remark}

Near-field energy can move back and forth around the source without escaping. The
radiative part is singled out by the limiting procedure $r\to\infty$: only energy that
still crosses arbitrarily large spheres is energy that has truly left the source.

% ==============================================================================
\subsection{From Radiation Energy to Field-Theory Energy}
\label{subsec:lec22:field-theory-motivation}
% ==============================================================================

\subsubsection*{Conceptual Background}
The Poynting vector is familiar, but the course has developed electrodynamics from an
action principle. We therefore want to derive energy and momentum from the same
language. The general strategy mirrors analytical mechanics:
\begin{itemize}
  \item absence of explicit time dependence gives energy conservation;
  \item absence of explicit spatial dependence gives momentum conservation;
  \item in relativistic field theory these combine into one tensorial conservation law.
\end{itemize}

The central object is the energy-momentum tensor. The lecture first derives it for a
general classical field theory and only later specializes it to the electromagnetic
field.

% ==============================================================================
\subsection{General Classical Field Theory}
\label{subsec:lec22:general-action}
% ==============================================================================

\subsubsection*{Mathematical Setup}
Let $\psi(x)$ be a classical field. To keep the derivation uncluttered, first assume
that $\psi$ has no internal index. The electromagnetic case is recovered by replacing
$\psi$ with the four fields $A_\mu$ and summing over that additional index.

We take the action to be
\begin{equation}\label{eq:lec22:general-action}
  \boxed{
  S[\psi]=\frac{1}{c}\int d^4x\,\mathcal{L}
  \qty(\psi,\partial_\mu\psi)}
\end{equation}
with no explicit dependence on $x^\mu$:
\begin{equation}\label{eq:lec22:no-explicit-x}
  \pdv{\mathcal{L}}{x^\mu}\bigg|_{\psi,\partial\psi}=0.
\end{equation}
\textit{Significance:} Equation \eqref{eq:lec22:no-explicit-x} encodes spacetime
homogeneity. Moving the whole closed system in space or time does not change the
physics.

The Lagrangian density depends only on the field and its first derivatives. Including
higher derivatives would usually produce higher-order equations of motion and require
extra boundary data.

% ==============================================================================
\subsection{Euler--Lagrange Equation for a Field}
\label{subsec:lec22:euler-lagrange}
% ==============================================================================

\subsubsection*{Derivation}
Vary the field,
\begin{equation}\label{eq:lec22:field-variation}
  \psi(x)\to \psi(x)+\delta\psi(x),
\end{equation}
with $\delta\psi=0$ on the boundary of the spacetime integration region. To first
order,
\begin{align}
  \delta S
  &= \frac{1}{c}\int d^4x\,
     \qty[
       \pdv{\mathcal{L}}{\psi}\,\delta\psi
       + \pdv{\mathcal{L}}{(\partial_\mu\psi)}\,\delta(\partial_\mu\psi)
     ] \notag\\
  &= \frac{1}{c}\int d^4x\,
     \qty[
       \pdv{\mathcal{L}}{\psi}\,\delta\psi
       + \pdv{\mathcal{L}}{(\partial_\mu\psi)}\,\partial_\mu(\delta\psi)
     ].
  \label{eq:lec22:variation-before-parts}
\end{align}
Now integrate the second term by parts:
\begin{align}
  \pdv{\mathcal{L}}{(\partial_\mu\psi)}\,\partial_\mu(\delta\psi)
  &=
  \partial_\mu\qty[
    \pdv{\mathcal{L}}{(\partial_\mu\psi)}\,\delta\psi
  ]
  -
  \partial_\mu\qty[
    \pdv{\mathcal{L}}{(\partial_\mu\psi)}
  ]\delta\psi.
  \label{eq:lec22:parts-identity}
\end{align}
The total derivative becomes a boundary term and vanishes because $\delta\psi$ is fixed
to be zero on the boundary. Therefore
\begin{equation}\label{eq:lec22:variation-after-parts}
  \delta S
  = \frac{1}{c}\int d^4x\,
    \qty[
      \pdv{\mathcal{L}}{\psi}
      - \partial_\mu\qty(\pdv{\mathcal{L}}{(\partial_\mu\psi)})
    ]\delta\psi.
\end{equation}
Since the action must be stationary for every allowed $\delta\psi$,
\begin{equation}\label{eq:lec22:field-euler-lagrange}
  \boxed{
  \partial_\mu\qty(\pdv{\mathcal{L}}{(\partial_\mu\psi)})
  =
  \pdv{\mathcal{L}}{\psi}.}
\end{equation}
\textit{Significance:} This is the field-theory Euler--Lagrange equation. For several
fields $\psi_a$, the same equation holds for each label $a$.

% ==============================================================================
\subsection{Translation Symmetry and the Canonical Tensor}
\label{subsec:lec22:canonical-tensor}
% ==============================================================================

\subsubsection*{Mathematical Setup}
Although $\mathcal{L}$ has no explicit $x^\mu$ dependence, it still changes with
$x^\mu$ through the field configuration $\psi(x)$. The total derivative is therefore
\begin{equation}\label{eq:lec22:L-chain-rule}
  \partial_\mu\mathcal{L}
  =
  \pdv{\mathcal{L}}{\psi}\,\partial_\mu\psi
  +
  \pdv{\mathcal{L}}{(\partial_\nu\psi)}\,\partial_\mu\partial_\nu\psi.
\end{equation}
Use the Euler--Lagrange equation \eqref{eq:lec22:field-euler-lagrange} to replace
$\partial\mathcal{L}/\partial\psi$:
\begin{align}
  \partial_\mu\mathcal{L}
  &=
  \partial_\nu\qty(\pdv{\mathcal{L}}{(\partial_\nu\psi)})\partial_\mu\psi
  +
  \pdv{\mathcal{L}}{(\partial_\nu\psi)}\,\partial_\nu\partial_\mu\psi \notag\\
  &=
  \partial_\nu\qty[
    \pdv{\mathcal{L}}{(\partial_\nu\psi)}\,\partial_\mu\psi
  ].
  \label{eq:lec22:L-total-divergence}
\end{align}
On the other hand,
\begin{equation}\label{eq:lec22:L-delta}
  \partial_\mu\mathcal{L}
  = \partial_\nu\qty(\delta^\nu{}_\mu\,\mathcal{L}).
\end{equation}
Equating \eqref{eq:lec22:L-total-divergence} and \eqref{eq:lec22:L-delta} gives
\begin{equation}\label{eq:lec22:canonical-conservation-step}
  \partial_\nu\qty[
    \pdv{\mathcal{L}}{(\partial_\nu\psi)}\,\partial_\mu\psi
    -\delta^\nu{}_\mu\,\mathcal{L}
  ]=0.
\end{equation}

\dfn[def:lec22:canonical-T]{Canonical Energy-Momentum Tensor}{
For a field $\psi$ with Lagrangian density $\mathcal{L}(\psi,\partial_\mu\psi)$,
\begin{equation}\label{eq:lec22:canonical-T}
  \boxed{
  T^\nu{}_\mu
  =
  \pdv{\mathcal{L}}{(\partial_\nu\psi)}\,\partial_\mu\psi
  -\delta^\nu{}_\mu\,\mathcal{L}.}
\end{equation}
It obeys the local conservation law
\begin{equation}\label{eq:lec22:T-conservation}
  \boxed{\partial_\nu T^\nu{}_\mu=0.}
\end{equation}
}

\textit{Significance:} For each fixed $\mu$, \eqref{eq:lec22:T-conservation} is a
continuity equation. The four values of $\mu$ correspond to conservation of energy and
the three components of momentum.

\nt{The object in \eqref{eq:lec22:canonical-T} is the canonical tensor. In gauge
theories, including electromagnetism, it is often improved or symmetrized before being
identified with the physical stress-energy tensor. The lecture's point here is the
Noether origin: translation symmetry produces a local conserved current with two
spacetime indices.}

% ==============================================================================
\subsection*{Summary}

\begin{itemize}
  \item The radiation fields of an electric dipole are transverse:
        $\bE\perp\bB\perp\uvec{r}$ and $|\bE|=|\bB|$.
  \item The Poynting vector for dipole radiation is radial and scales as $1/r^2$:
        \[
          \bvec{S}
          = \frac{|\ddot{\bvec{p}}(t_r)|^2\sin^2\theta}{4\pi c^3 r^2}\,\uvec{r}.
        \]
  \item The angular distribution is proportional to $\sin^2\theta$, so no radiation is
        emitted along the dipole axis.
  \item The total power radiated to infinity is
        \[
          P=\frac{2}{3c^3}|\ddot{\bvec{p}}(t_r)|^2.
        \]
  \item Radiation does not violate energy conservation; the energy comes from work done
        by whatever force accelerates the charges.
  \item A general first-derivative field theory with
        $S=(1/c)\int d^4x\,\mathcal{L}(\psi,\partial_\mu\psi)$ gives the field
        Euler--Lagrange equation
        \[
          \partial_\mu\qty(\pdv{\mathcal{L}}{(\partial_\mu\psi)})
          =\pdv{\mathcal{L}}{\psi}.
        \]
  \item If $\mathcal{L}$ has no explicit $x^\mu$ dependence, translation symmetry gives
        the canonical energy-momentum tensor
        \[
          T^\nu{}_\mu
          =
          \pdv{\mathcal{L}}{(\partial_\nu\psi)}\,\partial_\mu\psi
          -\delta^\nu{}_\mu\,\mathcal{L},
          \qquad
          \partial_\nu T^\nu{}_\mu=0.
        \]
\end{itemize}
